Las startups que desarrollan servicios sobre Node.js rara vez cuentan con un equipo de seguridad dedicado. Este trabajo presenta logSguarDian, una librería npm que implementa un middleware de autoprotección en tiempo de ejecución (RASP) para Express y detecta inyección SQL, \textit{Cross-Site Scripting}, \textit{Path Traversal} e inyección de comandos. El middleware extrae 76 características de cada petición y las evalúa en \texttt{worker\_threads} con un modelo híbrido. Un Random Forest decide el bloqueo y un Isolation Forest registra anomalías sin bloquear. Se entrenaron con 268\,240 de las 383\,202 filas del corpus y se evaluaron en un conjunto de prueba bloqueado y en cinco aplicaciones.

El extractor requiere 0.0273 ms por petición (p95; límite de 5 ms). El Random Forest alcanza un F1 macro de 0.9776 y el Isolation Forest, un recall de 0.9125 con 0.0546 de falsos positivos. La cobertura por categoría supera el 80\,\% en inyección SQL, \textit{Path Traversal} e inyección de comandos con el modelo evaluado. La latencia cumple el criterio de 15 ms en carga normal y de ataque, pero no en volumen, y no se cumplen los criterios de memoria y CPU.

Con un \textit{User-Agent} de navegador, el modelo validado marcó para bloqueo el 97.1\,\% de las peticiones benignas en Juice Shop (modo \textit{monitor}) y un modelo no documentado bloqueó el 92.5\,\% en dvws-node; el corpus benigno es la debilidad principal. Las particiones comparten casi-duplicados y el conjunto de prueba se leyó tres veces; las métricas miden separación dentro del corpus.
